% 65e Olympiade Internationale de Mathématiques — Bath, Royaume-Uni.
% Énoncés officiels en français (imo-official.org). Chaque problème vaut 7 points.

\begin{center}
  \emph{Premier jour --- mardi 16 juillet 2024 --- Durée : 4 heures 30}
\end{center}

\subsection*{Problème 1}
Déterminer tous les réels $\alpha$ tels que, pour tout entier $n$ strictement positif, l'entier
\[
\lfloor \alpha \rfloor + \lfloor 2\alpha \rfloor + \cdots + \lfloor n\alpha \rfloor
\]
soit un multiple de $n$. (On rappelle que $\lfloor z \rfloor$ désigne le plus grand entier inférieur ou égal
à $z$. Par exemple, $\lfloor -\pi \rfloor = -4$ et $\lfloor 2 \rfloor = \lfloor 2{,}9 \rfloor = 2$.)

\subsection*{Problème 2}
Déterminer tous les couples $(a, b)$ d'entiers strictement positifs pour lesquels il existe des entiers
strictement positifs $g$ et $N$ tels que l'égalité
\[
\operatorname{PGCD}(a^n + b, b^n + a) = g
\]
soit vérifiée pour tout entier $n \geqslant N$. (On rappelle que $\operatorname{PGCD}(x, y)$ désigne le
plus grand commun diviseur des entiers $x$ et $y$.)

\subsection*{Problème 3}
Soit $a_1, a_2, a_3, \ldots$ une suite d'entiers strictement positifs, et soit $N$ un entier strictement
positif. On suppose que, pour tout $n > N$, $a_n$ est égal au nombre de fois que la valeur $a_{n-1}$
apparaît dans la liste $a_1, a_2, \ldots, a_{n-1}$.

Démontrer que parmi les deux suites $a_1, a_3, a_5, \ldots$ et $a_2, a_4, a_6, \ldots$, au moins une est
ultimement périodique.

(Une suite $b_1, b_2, b_3, \ldots$ est \emph{ultimement périodique} s'il existe des entiers strictement
positifs $p$ et $M$ tels que $b_{m+p} = b_m$ pour tout entier $m \geqslant M$.)

\bigskip
\begin{center}
  \emph{Second jour --- mercredi 17 juillet 2024 --- Durée : 4 heures 30}
\end{center}

\subsection*{Problème 4}
Soit $ABC$ un triangle vérifiant $AB < AC < BC$. On note $\omega$ le cercle inscrit au triangle $ABC$, et
$I$ le centre de $\omega$. Soit $X$ le point de la droite $(BC)$, distinct de $C$, telle que la parallèle à
$(AC)$ passant par $X$ soit tangente à $\omega$. Similairement, soit $Y$ le point de la droite $(BC)$,
distinct de $B$, telle que la parallèle à $(AB)$ passant par $Y$ soit tangente à $\omega$. La droite $(AI)$
recoupe le cercle circonscrit au triangle $ABC$ en $P$ (distinct de $A$). Soit $K$ et $L$ les milieux des
segments $[AC]$ et $[AB]$ respectivement.

Démontrer que $\widehat{KIL} + \widehat{YPX} = 180^\circ$.

\subsection*{Problème 5}
Turbo l'escargot joue à un jeu sur une grille avec $2024$ lignes et $2023$ colonnes. Il y a des monstres
cachés dans $2022$ cases. Initialement, Turbo ne connaît l'emplacement d'aucun monstre, mais il sait qu'il
y a exactement un monstre dans chaque ligne à part dans la première ligne et dans la dernière ligne, et que
chaque colonne contient au plus un monstre.

Turbo enchaîne plusieurs tentatives pour aller de la première ligne à la dernière ligne. À chaque tentative,
il choisit la case de la première ligne sur laquelle il commence, et se déplace de cases en cases, en allant
à chaque fois d'une case à une case ayant un côté en commun avec celle-ci. (Il a le droit de repasser sur une
case déjà visitée.) S'il accède à une case contenant un monstre, sa tentative se termine et il est transporté
sur la première ligne pour commencer une nouvelle tentative. Les monstres ne bougent pas, et Turbo se
souvient si chaque case qu'il a visitée contient un monstre ou non. S'il atteint n'importe quelle case de la
dernière ligne, sa tentative se termine et le jeu est fini.

Déterminer la valeur minimale de $n$ pour laquelle Turbo a une stratégie qui lui assure d'atteindre la
dernière ligne à sa $n$\textsuperscript{e} tentative ou avant, quelle que soit de la localisation des monstres.

\subsection*{Problème 6}
On note $\mathbb{Q}$ l'ensemble des nombres rationnels. Une fonction $f \colon \mathbb{Q} \to \mathbb{Q}$ est
dite \emph{aquesuliane} si elle vérifie la propriété suivante : pour tout $x, y \in \mathbb{Q}$, au moins
une des deux égalités est vérifiée :
\[
f(x + f(y)) = f(x) + y \qquad \text{ou} \qquad f(f(x) + y) = x + f(y).
\]
Démontrer qu'il existe un entier $c$ tel que pour toute fonction aquesuliane $f$, la quantité
$f(r) + f(-r)$ prenne au plus $c$ valeurs distinctes lorsque $r$ parcourt $\mathbb{Q}$, et déterminer la
plus petite valeur possible pour $c$ vérifiant cela.
